\documentclass[11pt,a4paper]{article}
\usepackage{amsmath}\usepackage{amssymb}\usepackage[table]{xcolor}
\usepackage{tabularx}\usepackage{multirow}\usepackage{geometry}
\usepackage{enumitem}\usepackage{parskip}\usepackage{titlesec}
\usepackage{tocloft}\usepackage[most]{tcolorbox}
\usepackage{tikz}\usepackage{pgfplots}\pgfplotsset{compat=1.18}\usetikzlibrary{arrows.meta,calc,positioning,shapes.geometric,patterns}
\geometry{margin=2.5cm}
\titleformat{\section}{\large\bfseries}{}{0em}{}
\setlength{\cftbeforesecskip}{6pt}
\definecolor{q1color}{HTML}{D6EAF8}\definecolor{q2color}{HTML}{D5F5E3}
\definecolor{q3color}{HTML}{FDEBD0}\definecolor{q4color}{HTML}{F9E79F}
\definecolor{q5color}{HTML}{E8DAEF}\definecolor{q6color}{HTML}{FADBD8}
\definecolor{q7color}{HTML}{D1F2EB}\definecolor{q8color}{HTML}{FCF3CF}
\definecolor{q9color}{HTML}{EBDEF0}\definecolor{q10color}{HTML}{F5EEF8}
\newenvironment{mstab}[1]{%
  {\color{#1!75!black}\nobreak Mark Scheme}\par\nopagebreak\smallskip
  \arrayrulecolor{#1!70!black}\renewcommand{\arraystretch}{1.3}%
  \tabularx{\textwidth}{|>{\columncolor{#1!12}\raggedright\arraybackslash}p{1.9cm}|>{\columncolor{#1!12}}X|>{\columncolor{#1!12}\centering\arraybackslash}p{1.1cm}|}%
  \hline\rowcolor{#1!45} Question & Answer & Marks \\ \hline}{%
  \endtabularx\arrayrulecolor{black}\par\medskip}
\newtcolorbox{ansbox}[1]{colback=#1!8,colbacktitle=#1!30,colframe=#1!70!black,coltitle=black,fonttitle=\normalfont,title=Worked Solution,boxrule=0.6pt,arc=2mm,left=4pt,right=4pt,top=4pt,bottom=4pt,breakable}
\newcommand{\topiclabel}[1]{\hfill\normalfont\small\itshape\textcolor{black!55}{#1}}
\newcommand{\examq}[4]{\section[#1 - Q#2]{#1 - Q#2 - #3 - #4}}
\begin{document}
\begin{center}
  {\LARGE \textbf{Cambridge O Level Physics 5054/12/M/J/16}}\\[6pt]
  {\large May/June 2016 -- Multiple Choice: Answers \& Solutions}
\end{center}
\bigskip

%=================================================================
\examq{5054/12/M/J/16}{1}{Physical Quantities and Measurements}{1}

\begin{mstab}{q1color}
1 & D & 1 \\ \hline
\end{mstab}

\begin{ansbox}{q1color}
The circumference is the distance \emph{all the way round} the curved surface of the ball, so the instrument must be able to follow that curve and read the length directly.

A tape is flexible: it can be wrapped round the ball and the circumference read straight off its scale. A rule is rigid and can measure only a straight line, so at best it gives the diameter. Calipers and a micrometer are designed to grip a body across its width, so they too measure the \emph{diameter}.

(A diameter can of course be turned into a circumference with $C = \pi d$, but that is a calculation, not a reading, and the question asks which instrument \emph{gives a reading that is} the circumference.)
\end{ansbox}

%=================================================================
\examq{5054/12/M/J/16}{2}{Forces or Dynamics}{1}

\begin{mstab}{q2color}
2 & B & 1 \\ \hline
\end{mstab}

\begin{ansbox}{q2color}
First find the \emph{resultant} force on the sledge. The man pulls forwards with $60\,\mathrm{N}$ and friction opposes the motion with $20\,\mathrm{N}$, so
\[ F = 60 - 20 = 40\,\mathrm{N}. \]
Now apply Newton's second law, $F = ma$:
\[ a = \frac{F}{m} = \frac{40}{25} = 1.6\,\mathrm{m/s^2}. \]
(Using the $60\,\mathrm{N}$ pull on its own, without subtracting the friction, gives the wrong answer $2.4\,\mathrm{m/s^2}$; dividing $25$ by $40$ instead gives $0.63\,\mathrm{m/s^2}$. Both are offered as distractors.)
\end{ansbox}

%=================================================================
\examq{5054/12/M/J/16}{3}{Forces or Dynamics}{1}

\begin{mstab}{q3color}
3 & B & 1 \\ \hline
\end{mstab}

\begin{ansbox}{q3color}
A body moving in a circle is continually changing the \emph{direction} of its velocity, so it is accelerating even though its speed is constant. By $F = ma$ the resultant force points the same way as that acceleration.

For circular motion the acceleration -- and hence the resultant force -- is directed towards the \emph{centre} of the circle. It is the sideways (centripetal) push of the road on the tyres that provides it.

So the arrow required is the one pointing from the car towards the centre of the circle turned by the car, which is arrow B. An arrow along the direction of motion (tangential) would speed the car up, and an arrow pointing outwards, away from the centre, would push the car off the bend -- neither is the resultant force here.
\end{ansbox}

%=================================================================
\examq{5054/12/M/J/16}{4}{Mass Weight and Density \textperiodcentered\ Forces or Dynamics}{1}

\begin{mstab}{q4color}
4 & B & 1 \\ \hline
\end{mstab}

\begin{ansbox}{q4color}
The property that resists a change in motion is \textbf{inertia}, and the measure of a body's inertia is its \textbf{mass}. This is Newton's first law read the other way round: the larger the mass, the larger the force needed to produce a given acceleration, since $a = F/m$.

The others are not measures of inertia. Weight is the gravitational \emph{force} on the body, so it changes with the field strength $g$ while the body's reluctance to accelerate does not. Volume is the space the body occupies, and density is mass per unit volume -- a small dense body and a large light body can have the same density but very different inertias.
\end{ansbox}

%=================================================================
\examq{5054/12/M/J/16}{5}{Mass Weight and Density}{1}

\begin{mstab}{q5color}
5 & C & 1 \\ \hline
\end{mstab}

\begin{ansbox}{q5color}
A newton meter is a spring balance. It is stretched by the gravitational pull on the brick, so what it responds to -- and what its scale is marked in -- is \emph{force}, measured in newtons. It therefore measures \textbf{weight}.

A beam balance compares the brick against known masses on the other arm. Because both sides are pulled by the same gravitational field, the comparison is unaffected by the value of $g$, and the reading is in kilograms. It therefore measures \textbf{mass}.

So: newton meter $\rightarrow$ weight, beam balance $\rightarrow$ mass.
\end{ansbox}

%=================================================================
\examq{5054/12/M/J/16}{6}{Turning Effect of Forces}{1}

\begin{mstab}{q6color}
6 & B & 1 \\ \hline
\end{mstab}

\begin{ansbox}{q6color}
The plank is uniform and pivoted at its \emph{mid-point}, so its own weight acts at the pivot and has no moment about it. Only the two added weights and the force at X need be considered.

Take moments about the pivot. The $12\,\mathrm{N}$ weight acts $2.0\,\mathrm{m}$ from the pivot on one side and the $8.0\,\mathrm{N}$ weight acts $2.0\,\mathrm{m}$ from the pivot on the other side:
\[ 12 \times 2.0 = 24\,\mathrm{N\,m}, \qquad 8.0 \times 2.0 = 16\,\mathrm{N\,m}. \]
The $12\,\mathrm{N}$ side therefore wins by
\[ 24 - 16 = 8.0\,\mathrm{N\,m}, \]
so that side of the plank tips downwards. Point X lies $4.0\,\mathrm{m}$ from the pivot on the \emph{same} side as the $12\,\mathrm{N}$ weight, so to cancel that surplus moment the force there must act \textbf{upwards}:
\[ F \times 4.0 = 8.0 \quad\Rightarrow\quad F = \frac{8.0}{4.0} = 2.0\,\mathrm{N}. \]
A force of $2.0\,\mathrm{N}$ upwards. (Getting the direction the wrong way round gives option A; dividing the $8.0\,\mathrm{N\,m}$ by $2.0\,\mathrm{m}$ instead of $4.0\,\mathrm{m}$ gives the $4.0\,\mathrm{N}$ options.)
\end{ansbox}

%=================================================================
\examq{5054/12/M/J/16}{7}{Physical Quantities and Measurements}{1}

\begin{mstab}{q7color}
7 & B & 1 \\ \hline
\end{mstab}

\begin{ansbox}{q7color}
A vector has both magnitude \emph{and} direction; a scalar has magnitude only.

Energy of every kind is a scalar -- there is no such thing as energy pointing north. Speed is a scalar too: it is the magnitude of the velocity, with the direction thrown away (velocity is the vector, speed is not). Time is a scalar, and in any case has no direction in space.

Force is the odd one out: a $10\,\mathrm{N}$ push to the left and a $10\,\mathrm{N}$ push to the right have utterly different effects, so a force is not fully specified until its direction is given. Force is a \textbf{vector}.
\end{ansbox}

%=================================================================
\examq{5054/12/M/J/16}{8}{Work Energy and Power \textperiodcentered\ The Nuclear Atom \textperiodcentered\ The Stars \& Universe}{1}

\begin{mstab}{q8color}
8 & C & 1 \\ \hline
\end{mstab}

\begin{ansbox}{q8color}
The Sun's core is at an enormous temperature and pressure, so hydrogen nuclei are forced close enough together to join and form helium. Joining light nuclei into a heavier one is \textbf{nuclear fusion}, and the small loss of mass in the process is released as energy.

The alternatives fail. Burning is a chemical reaction needing oxygen, and could not release anything like the Sun's output. Nuclear fission is the \emph{splitting} of heavy nuclei such as uranium, which is what happens in a nuclear power station on Earth, not in a star. And radiation is how the energy \emph{travels} to us once it has been produced, not the process that produces it.
\end{ansbox}

%=================================================================
\examq{5054/12/M/J/16}{9}{Work Energy and Power}{1}

\begin{mstab}{q9color}
9 & C & 1 \\ \hline
\end{mstab}

\begin{ansbox}{q9color}
Work done $=$ force $\times$ distance moved \emph{in the direction of the force}. The $300\,\mathrm{N}$ force is applied along the ramp, and from X to Y the box moves $5.0\,\mathrm{m}$ along that ramp, so those two are already in the same direction:
\[ W = F \times d = 300 \times 5.0 = 1500\,\mathrm{J}. \]
The $4.0\,\mathrm{m}$ and $3.0\,\mathrm{m}$ are the vertical and horizontal sides of the triangle. They are there to show that the slope length is $5.0\,\mathrm{m}$ (by Pythagoras) and to tempt you into using the wrong distance -- $300 \times 4.0 = 1200\,\mathrm{J}$ and $300 \times 3.0 = 900\,\mathrm{J}$ are both offered as wrong options. The vertical height would be used only with the \emph{weight} of the box, which is not given.
\end{ansbox}

%=================================================================
\examq{5054/12/M/J/16}{10}{Pressure}{1}

\begin{mstab}{q1color}
10 & C & 1 \\ \hline
\end{mstab}

\begin{ansbox}{q1color}
The pressure due to a liquid is
\[ p = \rho g h, \]
where $h$ is the \emph{depth below the free surface}. All four beakers hold the same liquid, so $\rho$ is the same and $g$ is the same; the pressure therefore depends on the depth alone.

It does not depend on the width or shape of the beaker, nor on how much liquid is present -- only on how far down the point lies. The greatest pressure is therefore at the point with the greatest depth below the surface of its own liquid, which is point C.
\end{ansbox}

%=================================================================
\examq{5054/12/M/J/16}{11}{Pressure}{1}

\begin{mstab}{q2color}
11 & B & 1 \\ \hline
\end{mstab}

\begin{ansbox}{q2color}
Pressure is defined as force per unit area:
\[ p = \frac{F}{A}. \]
The force pressing down on the ground is the block's \emph{weight}, and the weight $W$ is given directly, so no conversion is needed:
\[ p = \frac{W}{A}. \]
The gravitational field strength $g$ is a red herring. It would be needed only if the \emph{mass} had been given, since then $F = mg$; here the weight is already a force in newtons, so dividing by $g$ (options C and D) would be double-counting. Multiplying by the area instead of dividing (option A) gives the wrong units altogether.
\end{ansbox}

%=================================================================
\examq{5054/12/M/J/16}{12}{Kinetic Model of Matter \textperiodcentered\ Pressure}{1}

\begin{mstab}{q3color}
12 & A & 1 \\ \hline
\end{mstab}

\begin{ansbox}{q3color}
Work in units of ``atmospheres'', using the information that $10\,\mathrm{m}$ of water gives a pressure equal to one atmosphere.

At the \emph{surface} the bubble feels atmospheric pressure only:
\[ p_1 = 1\ \text{atm}, \qquad V_1 = 6.0\,\mathrm{cm^3}. \]
At the \emph{bottom}, $20\,\mathrm{m}$ down, the bubble feels atmospheric pressure \emph{plus} the pressure of $20\,\mathrm{m}$ of water, and $20\,\mathrm{m}$ of water is worth $2$ atmospheres:
\[ p_2 = 1 + 2 = 3\ \text{atm}. \]
The temperature is unchanged, so Boyle's law applies:
\[ p_1V_1 = p_2V_2 \quad\Rightarrow\quad V_2 = \frac{p_1V_1}{p_2} = \frac{1 \times 6.0}{3} = 2.0\,\mathrm{cm^3}. \]
The bubble is squeezed smallest where the pressure is greatest, at the bottom -- so the answer must be less than $6.0\,\mathrm{cm^3}$, which rules out the two larger options immediately.
\end{ansbox}

%=================================================================
\examq{5054/12/M/J/16}{13}{Kinetic Model of Matter}{1}

\begin{mstab}{q4color}
13 & C & 1 \\ \hline
\end{mstab}

\begin{ansbox}{q4color}
Heating a gas increases the average kinetic energy of its molecules, and since $E_k = \tfrac{1}{2}mv^2$ with the molecular mass fixed, that means the molecules \textbf{move faster}.

The other options misdescribe what happens. Moving faster means the molecules reach the walls more often, so they collide \emph{more} frequently, not less (which is part of why the pressure rises). Individual molecules do not expand -- it is the gas as a whole that would expand if it could. And the container has a \emph{fixed volume}, so the molecules cannot spread out: their average separation is unchanged.
\end{ansbox}

%=================================================================
\examq{5054/12/M/J/16}{14}{Thermal Properties \& Temperature}{1}

\begin{mstab}{q5color}
14 & C & 1 \\ \hline
\end{mstab}

\begin{ansbox}{q5color}
The two thermometers are identical apart from the bore of the capillary tube, so for a given temperature rise the \emph{volume} of liquid pushed out of the bulb is the same in each. What differs is how far along the tube that volume reaches:
\[ \Delta L = \frac{\Delta V}{A}, \]
where $A$ is the cross-sectional area of the bore.

\textbf{Sensitivity} is the length of thread produced per degree, $\Delta L / \Delta \theta$. The \emph{narrower} bore gives the longer thread for the same $\Delta V$, so the more sensitive thermometer is \textbf{P}.

\textbf{Range} is the number of degrees that fit into the finite length of tube available. Because P uses up its tube quickly, it runs out of scale sooner; Q, with its wider bore, advances only a little per degree and so can cover many more degrees before the thread reaches the top. The larger range belongs to \textbf{Q}.

Larger range Q, greater sensitivity P. (The two properties always pull against each other in a liquid-in-glass thermometer: you buy sensitivity at the cost of range.)
\end{ansbox}

%=================================================================
\examq{5054/12/M/J/16}{15}{Transfer of Thermal Energy}{1}

\begin{mstab}{q6color}
15 & D & 1 \\ \hline
\end{mstab}

\begin{ansbox}{q6color}
Take the three processes in turn.

\textbf{Conduction} needs the two bodies to be in contact through a material that conducts. The bread is not touching the grill; between them is air, which is a very poor conductor. So conduction delivers almost nothing.

\textbf{Convection} carries energy by the bulk movement of a heated fluid, and hot air \emph{rises}. The grill is above the bread, so any convection current it sets up moves upwards, away from the bread rather than towards it.

\textbf{Radiation} needs no medium at all and travels in straight lines in every direction, including downwards. A red-hot grill is a strong emitter of infra-red, and that radiation crosses the gap directly to the bread, which absorbs it and heats up.

So the energy reaches the bread by \textbf{radiation only}.
\end{ansbox}

%=================================================================
\examq{5054/12/M/J/16}{16}{Transfer of Thermal Energy \textperiodcentered\ Kinetic Model of Matter}{1}

\begin{mstab}{q7color}
16 & D & 1 \\ \hline
\end{mstab}

\begin{ansbox}{q7color}
The air near the cooling unit is \emph{cooled}. Cooling removes kinetic energy from the molecules, so they move more slowly and no longer push each other as far apart: on average they \textbf{move closer together}. The molecules themselves are unchanged -- particles never shrink or swell; it is the spacing between them that alters.

The same number of molecules now occupies a smaller volume, so by
\[ \rho = \frac{m}{V} \]
the density of that air \textbf{increases}.

Move closer together, density increases. (This is exactly why the cooling unit is placed at the top of a refrigerator: the dense cold air sinks, warmer air rises to take its place, and a convection current circulates the cooling round the whole cabinet.)
\end{ansbox}

%=================================================================
\examq{5054/12/M/J/16}{17}{General Wave Properties}{1}

\begin{mstab}{q8color}
17 & D & 1 \\ \hline
\end{mstab}

\begin{ansbox}{q8color}
The \emph{frequency} of a wave is fixed by whatever is generating it -- the dipper in the ripple tank -- so it cannot change when the wave crosses into a new depth of water. The same number of crests enters the shallow region each second as leaves the deep region.

The \emph{speed} of a water wave does depend on the depth: waves travel more slowly in shallow water. So $v$ falls.

Since $v = f\lambda$ with $f$ fixed, a smaller $v$ means a smaller $\lambda$: the wavefronts bunch closer together. That bunching, together with the change of direction at the boundary, is what refraction is.

So the properties that change are \textbf{speed and wavelength}. (The amplitude is not what refraction alters, and the frequency, as argued, never changes -- which rules out the other three rows.)
\end{ansbox}

%=================================================================
\examq{5054/12/M/J/16}{18}{Light (Reflection/Refraction/Lenses)}{1}

\begin{mstab}{q9color}
18 & A & 1 \\ \hline
\end{mstab}

\begin{ansbox}{q9color}
Angles in reflection are always measured between the ray and the \textbf{normal} -- the line drawn at right angles to the mirror surface at the point where the ray strikes it. The line PQ in the diagram is that perpendicular, so PQ is the \textbf{normal}. (It is not the reflected ray: the reflected ray leaves the mirror on the other side of the normal, at an angle to it.)

The $40^{\circ}$ is marked between the incident ray and PQ, so it is the angle of incidence measured from the normal. The law of reflection then gives
\[ \text{angle of reflection} = \text{angle of incidence} = 40^{\circ}. \]
Angle of reflection $40^{\circ}$, line PQ is the normal.

(The $50^{\circ}$ options are the complement, $90^{\circ} - 40^{\circ}$ -- what you get by measuring from the mirror surface instead of from the normal.)
\end{ansbox}

%=================================================================
\examq{5054/12/M/J/16}{19}{Light (Reflection/Refraction/Lenses)}{1}

\begin{mstab}{q1color}
19 & B & 1 \\ \hline
\end{mstab}

\begin{ansbox}{q1color}
A parallel beam travelling along the principal axis is brought to a focus by a converging lens at the \emph{principal focus}. The \textbf{focal length} is defined as the distance from the \emph{centre of the lens} to that principal focus, measured along the principal axis.

So the length required runs from the lens itself to the point where the rays meet -- length B. Lengths measured from somewhere other than the centre of the lens, or that run from the lens to a point beyond the focus, do not give the focal length; nor does the distance between the two principal foci, one on each side of the lens, which is twice the focal length.
\end{ansbox}

%=================================================================
\examq{5054/12/M/J/16}{20}{Light (Reflection/Refraction/Lenses)}{1}

\begin{mstab}{q2color}
20 & B & 1 \\ \hline
\end{mstab}

\begin{ansbox}{q2color}
Judge each lens by what it does to the beam passing through it, not by the shape drawn:
\begin{itemize}[nosep,leftmargin=*]
  \item a \textbf{converging} (convex) lens bends every ray \emph{towards} the principal axis, so the beam emerging is narrower than the beam entering -- it is brought together;
  \item a \textbf{diverging} (concave) lens bends every ray \emph{away} from the axis, so the emerging beam is wider than the beam entering -- it is spread out.
\end{itemize}
Tracing the rays from the ray-box through the three positions, the beam is squeezed together at position 1 and again at position 2, so both of those lenses are converging; at position 3 the rays are pushed apart, so that lens is diverging.

Converging, converging, diverging.
\end{ansbox}

%=================================================================
\examq{5054/12/M/J/16}{21}{Light (Reflection/Refraction/Lenses)}{1}

\begin{mstab}{q3color}
21 & A & 1 \\ \hline
\end{mstab}

\begin{ansbox}{q3color}
Snell's law gives the refractive index as
\[ n = \frac{\sin(\text{angle in air})}{\sin(\text{angle in glass})}, \]
and both angles must be measured from the \textbf{normal} to the surface, never from the surface itself.

In the diagram the normal is the line drawn perpendicular to the glass face at the point where the ray enters. Of the four marked angles, $w$ is the one between the incident ray and that normal, and $y$ is the one between the refracted ray inside the glass and the same normal; $x$ and $z$ are measured from the \emph{surface}, so each is the complement of the angle actually needed.

Hence
\[ n = \frac{\sin w}{\sin y}. \]
A useful check: light entering glass from air bends \emph{towards} the normal, so $y < w$ and the ratio $\sin w / \sin y$ is greater than $1$ -- as a refractive index must be. The inverted ratio $\sin y / \sin w$ would be less than $1$, which rules out option C at once.
\end{ansbox}

%=================================================================
\examq{5054/12/M/J/16}{22}{Sound \textperiodcentered\ General Wave Properties}{1}

\begin{mstab}{q4color}
22 & B & 1 \\ \hline
\end{mstab}

\begin{ansbox}{q4color}
Sound is a \emph{longitudinal} wave: the layers of air are displaced back and forth along the direction of travel, so the lines in the diagram bunch together at the compressions and spread apart at the rarefactions.

The wavelength is the distance between two \emph{successive} points that are in phase -- in practice, from the middle of one compression to the middle of the \textbf{next} compression, or from one rarefaction to the next. One complete cycle must be included, no more and no less.

Reading the marked distances against that rule, the one that spans exactly one compression-plus-rarefaction, and so ends on the matching point of the neighbouring compression, is distance B. The distance from a compression to the adjacent rarefaction is only \emph{half} a wavelength, and a distance reaching the second compression beyond the first would be two wavelengths -- both are offered as distractors.
\end{ansbox}

%=================================================================
\examq{5054/12/M/J/16}{23}{Simple Magnetism and Magnetic Fields}{1}

\begin{mstab}{q5color}
23 & C & 1 \\ \hline
\end{mstab}

\begin{ansbox}{q5color}
A plotting compass is a tiny magnet free to turn, so it lines up along the magnetic field at its position, with its \textbf{north} pole pointing the way the field points. By convention the field lines of a bar magnet leave the N pole, curve round outside the magnet, and enter the S pole. (The Earth's field is to be ignored, so the bar magnet's field is the only one acting.)

So the method is: at each of the four positions work out which way the field line runs there, then check whether that matches the way the needle in the given picture is pointing. Off the ends of the magnet the lines point straight out of the N pole and straight into the S pole; alongside the magnet they run back along its length from the N end towards the S end.

The only one of the four positions at which the field direction matches the orientation of the compass shown is position C.
\end{ansbox}

%=================================================================
\examq{5054/12/M/J/16}{24}{Electric Charge}{1}

\begin{mstab}{q6color}
24 & A & 1 \\ \hline
\end{mstab}

\begin{ansbox}{q6color}
In any solid, the only charges that can move are \textbf{electrons}, which are \emph{negative}. Positive charges sit in fixed nuclei and are never transferred by rubbing, so options C and D can be discarded straight away.

The cloth ends up positively charged, which means it has \emph{lost} electrons. Those electrons must have gone onto the polythene rod, leaving it equally negatively charged.

So negative charges are transferred \textbf{from the cloth to the polythene}. (Charge is conserved: the rod gains exactly as much negative charge as the cloth loses, which is why the two always end up with equal and opposite charges.)
\end{ansbox}

%=================================================================
\examq{5054/12/M/J/16}{25}{Electrical Quantities}{1}

\begin{mstab}{q7color}
25 & A & 1 \\ \hline
\end{mstab}

\begin{ansbox}{q7color}
Charge is current $\times$ time, so current is charge divided by time -- with the time in \emph{seconds}:
\[ t = 3.0 \times 60 = 180\,\mathrm{s} \]
\[ I = \frac{Q}{t} = \frac{45}{180} = 0.25\,\mathrm{A}. \]
(Forgetting to convert the minutes gives $45/3.0 = 15\,\mathrm{A}$, and multiplying instead of dividing gives $45 \times 3.0 = 135\,\mathrm{A}$ -- both are offered as distractors.)
\end{ansbox}

%=================================================================
\examq{5054/12/M/J/16}{26}{Electrical Quantities}{1}

\begin{mstab}{q8color}
26 & D & 1 \\ \hline
\end{mstab}

\begin{ansbox}{q8color}
The potential difference between two points is the work done (energy transferred) per unit charge passing between them:
\[ V = \frac{W}{Q} \quad\Rightarrow\quad 1\,\mathrm{V} = \frac{1\,\mathrm{J}}{1\,\mathrm{C}}. \]
So one volt is \textbf{one joule per coulomb}.

The others are the wrong combinations: one coulomb per ampere is a second (from $Q = It$), one coulomb per joule is the reciprocal of a volt, and one joule per ampere is a weber (a magnetic flux), none of which is a potential difference.
\end{ansbox}

%=================================================================
\examq{5054/12/M/J/16}{27}{Electric Circuits}{1}

\begin{mstab}{q9color}
27 & A & 1 \\ \hline
\end{mstab}

\begin{ansbox}{q9color}
Read the diagram first: the three cells are connected in \textbf{parallel}, all their positive terminals joined together and all their negative terminals joined together.

Cells add their e.m.f.s only when they are in \emph{series}. In parallel every cell sits across the same pair of terminals, so the e.m.f. of the battery is that of a single cell, $1.5\,\mathrm{V}$, not $4.5\,\mathrm{V}$. That disposes of option D, and of option C as well: the energy given to each coulomb \emph{is} the e.m.f., so each coulomb receives the same $1.5\,\mathrm{J}$ as it would from one cell. Option B is wrong because fixed cells in parallel give one fixed e.m.f., not a range -- a variable e.m.f. would need a potential divider.

What parallel connection does give is \emph{capacity}. The current drawn by the circuit is shared between the three cells, so each supplies only a third of it and each is used up a third as fast. The battery can therefore \textbf{supply a current for a longer time} than a single cell.
\end{ansbox}

%=================================================================
\examq{5054/12/M/J/16}{28}{Electric Circuits \textperiodcentered\ Electrical Quantities}{1}

\begin{mstab}{q1color}
28 & B & 1 \\ \hline
\end{mstab}

\begin{ansbox}{q1color}
The power of the lamp is
\[ P = VI, \]
so the two readings needed are the potential difference \emph{across the lamp} and the current \emph{through the lamp}. That fixes how each meter must be connected.

\begin{itemize}[nosep,leftmargin=*]
  \item The \textbf{ammeter} must carry the lamp's current, so it goes in \emph{series} with the lamp. It has a very low resistance, so it hardly affects the circuit.
  \item The \textbf{voltmeter} must have the lamp's p.d. across it, so it goes in \emph{parallel} with the lamp. It has a very high resistance, so it draws almost no current.
\end{itemize}
Connecting them the other way round would be useless and dangerous: a voltmeter in series would block the current almost completely, and an ammeter in parallel across the lamp would short it out.

The diagram that has the ammeter in the main circuit loop with the lamp and the voltmeter bridged across the lamp alone is B.
\end{ansbox}

%=================================================================
\examq{5054/12/M/J/16}{29}{Electric Circuits}{1}

\begin{mstab}{q2color}
29 & A & 1 \\ \hline
\end{mstab}

\begin{ansbox}{q2color}
A voltmeter reads the difference in potential between the two points its leads touch. Here it must read $3.0\,\mathrm{V}$, which is exactly \emph{half} the $6.0\,\mathrm{V}$ supplied by the battery.

The resistors are identical, so in any series chain the supply voltage is shared equally between them: each resistor in a chain of $n$ takes $6.0/n$ volts. To pick up half the supply, lead L must therefore be moved to the point that is \emph{halfway round the network in resistance} from K -- the point with exactly half the total resistance of the current-carrying path between it and K. Note also that no p.d. appears across a resistor carrying no current, and none appears along a plain connecting wire, so points joined by wire alone are at the same potential and cannot differ by $3.0\,\mathrm{V}$.

Working round the network on that basis, the point at the midpoint of the chain relative to K -- and so the one giving a reading of $3.0\,\mathrm{V}$ -- is point A.
\end{ansbox}

%=================================================================
\examq{5054/12/M/J/16}{30}{Practical Electricity / Electrical Safety}{1}

\begin{mstab}{q3color}
30 & A & 1 \\ \hline
\end{mstab}

\begin{ansbox}{q3color}
Check the diagram against the rules for wiring a three-pin plug, and notice that three of the four statements describe things that are \emph{correct}:

\begin{itemize}[nosep,leftmargin=*]
  \item The earth wire E \emph{should} go to the longest pin, the one that makes contact first -- so option B describes no error.
  \item The fuse F \emph{should} be in the live wire L, so that it cuts the live connection when it blows and the appliance is left isolated from the live supply. Option C is therefore a correct feature, not a fault.
  \item The fuse is a single component in series; which of its two ends the live wire reaches first makes no electrical difference, so option D is not an error either.
\end{itemize}

What is wrong is mechanical. The cable grip S must clamp the outer insulating cover C of the cable, so that any pull on the cable is taken by the tough outer sheath and not by the three thin wires at their terminals -- above all not by the live wire. In the diagram the grip has closed on the inner wires instead, with the cable cover stopping short of it.

The error is that \textbf{the cable cover C is not under the clip S}.
\end{ansbox}

%=================================================================
\examq{5054/12/M/J/16}{31}{Practical Electricity / Electrical Safety}{1}

\begin{mstab}{q4color}
31 & C & 1 \\ \hline
\end{mstab}

\begin{ansbox}{q4color}
A ``unit'' of electricity is one kilowatt-hour, so the power must be in kilowatts and the time in hours.
\[ P = 100\,\mathrm{W} = 0.100\,\mathrm{kW} \]
\[ t = 5\ \text{hours per day} \times 21\ \text{days} = 105\,\mathrm{h} \qquad (3\ \text{weeks} = 3 \times 7 = 21\ \text{days}) \]
\[ E = Pt = 0.100 \times 105 = 10.5\,\mathrm{kW\,h} \]
\[ \text{cost} = 10.5 \times \$0.24 = \$2.52. \]
(Using three days instead of three weeks gives \$0.36, and leaving the power in watts gives \$25.20 -- a factor of ten out. Both appear among the options.)
\end{ansbox}

%=================================================================
\examq{5054/12/M/J/16}{32}{Electro Magnetism / Electromagnetic Induction}{1}

\begin{mstab}{q5color}
32 & D & 1 \\ \hline
\end{mstab}

\begin{ansbox}{q5color}
The sides of the coil carry a current in a magnetic field, so each feels a force
\[ F = BIL, \]
and the pair of forces provides the turning effect. Deal with the two requirements separately.

\textbf{Reversing the rotation.} Fleming's left-hand rule shows that the direction of $F$ reverses if the field is reversed, and also if the current is reversed. Reversing \emph{both} at once reverses the force twice over, leaving it exactly as it was -- so option B changes nothing. Reversing just one of them is what turns the motor the other way.

\textbf{Making it faster.} A bigger turning effect needs a bigger force, and $F = BIL$ grows with the current $I$ (and with the number of turns, since each turn adds its own pair of forces). So the current must be \emph{increased}; decreasing it, as in option C, would slow the motor down.

Option A increases the current and the turns but changes nothing about direction. Only \textbf{reversing the magnetic field and increasing the current} does both jobs at once.
\end{ansbox}

%=================================================================
\examq{5054/12/M/J/16}{33}{Electro Magnetism / Electromagnetic Induction}{1}

\begin{mstab}{q6color}
33 & A & 1 \\ \hline
\end{mstab}

\begin{ansbox}{q6color}
Use Fleming's left-hand rule: first finger along the \textbf{field}, second finger along the \textbf{current}, thumb then gives the \textbf{force} (thrust).

The field between the poles runs from the N pole across to the S pole. The current runs along the wire from P to Q, in the direction of the arrow. These two are at right angles to each other, both lying in the plane of the page, so the force -- which is perpendicular to \emph{both} -- must point at right angles to the page.

Setting the left hand with the first finger along the field and the second finger along the current, the thumb points away from the reader: the force on PQ is \textbf{into the page}.

A useful check: the force in the motor effect is never along the field and never along the current, so ``towards the N-pole'' and ``towards the S-pole'' can both be ruled out before any rule is applied. That leaves only into or out of the page.
\end{ansbox}

%=================================================================
\examq{5054/12/M/J/16}{34}{Electro Magnetism / Electromagnetic Induction}{1}

\begin{mstab}{q7color}
34 & D & 1 \\ \hline
\end{mstab}

\begin{ansbox}{q7color}
For a given power to be delivered,
\[ P = VI \quad\Rightarrow\quad I = \frac{P}{V}, \]
so raising the transmission voltage \emph{lowers} the current in the cables.

The power wasted as heat in the cables themselves is
\[ P_{\text{loss}} = I^2R, \]
which depends on the \emph{square} of the current. Halving the current cuts the loss to a quarter; raising the voltage by a factor of a hundred cuts the loss by a factor of ten thousand.

So a high-voltage system gives a \textbf{low} current in the cables and a \textbf{low} power loss. (This is why step-up and step-down transformers are used at each end of the grid, and incidentally why thinner, cheaper cables will do.)
\end{ansbox}

%=================================================================
\examq{5054/12/M/J/16}{35}{Uses of Oscilloscope}{1}

\begin{mstab}{q8color}
35 & C & 1 \\ \hline
\end{mstab}

\begin{ansbox}{q8color}
Thermionic emission is the escape of \textbf{electrons} from the surface of a hot metal. Heating the filament gives the free electrons inside it enough kinetic energy to break out through the surface, so the filament must be at a \textbf{high} temperature -- a cold filament emits nothing, which rules out options B and D.

The particles emitted are electrons, so they carry a \emph{negative} charge. Unlike charges attract, so a plate that attracts them must be charged \textbf{positive}. (A negative plate would repel them straight back to the filament.)

Positive plate, high filament temperature. This is exactly the electron gun of a cathode-ray tube: a heated cathode to liberate the electrons, and a positive anode to pull them away and accelerate them into a beam.
\end{ansbox}

%=================================================================
\examq{5054/12/M/J/16}{36}{Electric Circuits}{1}

\begin{mstab}{q9color}
36 & B & 1 \\ \hline
\end{mstab}

\begin{ansbox}{q9color}
The $6.0\,\Omega$ resistor and the variable resistor $R$ are in series across the $6.0\,\mathrm{V}$ supply, so the supply voltage is shared between them in proportion to their resistances:
\[ V_1 = 6.0 \times \frac{6.0}{6.0 + R}, \qquad V_2 = 6.0 \times \frac{R}{6.0 + R}, \qquad V_1 + V_2 = 6.0\,\mathrm{V}. \]
With $R = 6.0\,\Omega$ the two shares are equal, $V_1 = V_2 = 3.0\,\mathrm{V}$, as stated.

Now increase $R$. The total resistance rises, so the current
\[ I = \frac{6.0}{6.0 + R} \]
falls. The fixed resistor's share is $V_1 = I \times 6.0$, and with $I$ smaller that reading \textbf{decreases}. Since the two readings must still add to $6.0\,\mathrm{V}$, the other reading \textbf{increases} -- and indeed $R$ now takes the larger share of the supply because it is the larger resistance.

$V_1$ decreases, $V_2$ increases. (They cannot both move the same way, so options A and D can be rejected on the $V_1 + V_2 = 6.0\,\mathrm{V}$ argument alone.)
\end{ansbox}

%=================================================================
\examq{5054/12/M/J/16}{37}{Radioactivity}{1}

\begin{mstab}{q1color}
37 & A & 1 \\ \hline
\end{mstab}

\begin{ansbox}{q1color}
An alpha-particle is a helium nucleus, $^{4}_{2}\mathrm{He}$: two protons and two neutrons. Emitting one therefore takes $4$ from the nucleon number and $2$ from the proton number:
\[ ^{241}_{\ 95}\mathrm{Am} \rightarrow\ ^{237}_{\ 93}\mathrm{X} + ^{4}_{2}\mathrm{He}. \]
Check both numbers balance: $241 = 237 + 4$ and $95 = 93 + 2$.

The proton number is what fixes the element, and $Z = 93$ is neptunium, so the product is $^{237}_{\ 93}\mathrm{Np}$.

(The wrong options come from the two standard slips: \emph{adding} $2$ to the proton number instead of subtracting it, which gives berkelium, $Z = 97$; and leaving the nucleon number at $241$ or taking only $2$ off it, which gives the mass-$245$ options.)
\end{ansbox}

%=================================================================
\examq{5054/12/M/J/16}{38}{Radioactivity}{1}

\begin{mstab}{q2color}
38 & A & 1 \\ \hline
\end{mstab}

\begin{ansbox}{q2color}
Ionisation happens when the radiation knocks electrons off the atoms it passes, and how strongly it does so depends on its charge and on how slowly it moves through the material.

\begin{itemize}[nosep,leftmargin=*]
  \item An \textbf{alpha-particle} is a helium nucleus with charge $+2e$ and a comparatively large mass, so it interacts very strongly and ionises heavily. It gives up all its energy within a few centimetres of air, which is why a sheet of paper stops it.
  \item A \textbf{beta-particle} is an electron, charge $-e$ and far lighter, so it ionises much less and penetrates a few millimetres of aluminium.
  \item A \textbf{gamma-ray} has no charge and no mass at all, so it interacts least of the three, ionises only weakly, and needs thick lead to absorb it.
\end{itemize}

So ionising effect runs \textbf{alpha, beta, gamma} from highest to lowest -- the exact reverse of the order of penetrating power, which is the easiest way to remember it.
\end{ansbox}

%=================================================================
\examq{5054/12/M/J/16}{39}{Radioactivity}{1}

\begin{mstab}{q3color}
39 & D & 1 \\ \hline
\end{mstab}

\begin{ansbox}{q3color}
The half-life of a radioactive isotope is the time taken for \textbf{half} of the undecayed nuclei present to decay -- equivalently, the time for the number of nuclei of that isotope, or for the activity, to fall to half its value. That is statement D.

The others are each wrong in a particular way:
\begin{itemize}[nosep,leftmargin=*]
  \item decay is a random process whose probability per nucleus never alters, so the half-life is a \emph{constant} for a given isotope -- it does not change as the sample decays;
  \item the nucleon number is a property of the individual nucleus and does not halve at all; what halves is the \emph{number} of such nuclei remaining;
  \item the number of nuclei never actually reaches zero -- each half-life removes half of what is left, so there is no ``total decay time'' to take half of.
\end{itemize}
\end{ansbox}

%=================================================================
\examq{5054/12/M/J/16}{40}{The Nuclear Atom}{1}

\begin{mstab}{q4color}
40 & B & 1 \\ \hline
\end{mstab}

\begin{ansbox}{q4color}
Isotopes of the same element have the \emph{same} number of protons -- the proton number is what makes it that element -- but \emph{different} numbers of neutrons, and so different nucleon numbers.

So the neutrons can be ignored at first: simply count the protons in each of the three nuclei. Two of the diagrams contain the same number of protons as each other while the third contains a different number, and those first two then turn out to hold different numbers of neutrons, which is what makes them isotopes rather than identical nuclei.

The pair with matching proton numbers is \textbf{P and R}. Q has a different proton number, so it is a different element altogether and cannot be an isotope of either -- which also rules out ``P, Q and R''.
\end{ansbox}

\end{document}
